\documentclass[11pt]{article}
\usepackage[margin=1in]{geometry}
\usepackage{amsmath,amssymb,amsthm}
\usepackage[hidelinks]{hyperref}
\newtheorem{theorem}{Theorem}
\newcommand{\R}{\mathbb R}
\newcommand{\J}{J_0}
\title{A two-term counterexample to Bessel--Fourier kernel positivity}
\author{Research proof and Lean formalization}
\date{October 27, 2026}
\begin{document}
\maketitle
\begin{abstract}
We refute Acosta Babb's positivity conjecture for a modified Fej\'er kernel
of Bessel--Fourier series. At the strictly interior parameters
$\alpha=2/5$ and $x=47/50$, the two-term kernel is less than $-1/500$.
Exact rational series bounds identify the genuine first two positive
Bessel zeros and certify the negative value. The result is formally proved
in Lean without numerical root-location assumptions.
\end{abstract}

Let $j_m$ be the $m$th positive zero of the Bessel function $\J$.
Acosta Babb defines the order-zero kernel
\begin{equation}\label{kernel}
 K_M^\alpha(x)=\sum_{m=1}^{M}\left(1-\frac{m}{M+1}\right)
 j_m^\alpha\J(j_mx)
\end{equation}
and conjectures that it is nonnegative whenever
$0\leq\alpha\leq1/2$, $M\geq1$, and $0\leq x\leq1$
\cite[Conjecture~1, p.~11]{Babb26}.
The same assertion appears in the earlier preprint
\cite[Conjecture~1, p.~10]{Babb24}.
The published discussion fixes Bessel order zero, which is the case used here.

\begin{theorem}\label{main}
For the kernel \eqref{kernel},
\[
 K_2^{2/5}(47/50)<-\frac1{500}.
\]
In particular, the conjectured nonnegativity fails with both parameters
strictly inside their stated intervals.
\end{theorem}

\paragraph{Series and uniform error bounds.}
Put $q=x^2/4$ and write
\[
 \J(x)=\sum_{k\geq0}\frac{(-1)^kq^k}{(k!)^2},\qquad
 H(q)=\sum_{k\geq0}\frac{(-1)^kq^k}{k!(k+1)!}.
\]
Normal convergence on compact sets justifies
\[\J'(x)=-(x/2)H(x^2/4).\]
Let $p(q)$ and $h(q)$ be the respective sums for $0\leq k<20$.
For $0\leq q\leq9$, the absolute ratio of successive terms in either tail,
starting at index $20$, is at most $9/21^2<1/2$. Therefore
\begin{equation}\label{tails}
 \begin{aligned}
 \left|\J(x)-p(x^2/4)\right|&\leq
 e:=\frac{2\cdot9^{20}}{(20!)^2}<10^{-17},\\
 |H(q)-h(q)|&\leq
 e_H:=\frac{2\cdot9^{20}}{20!\,21!}<10^{-18}.
 \end{aligned}
\end{equation}
The first bound applies for $0\leq x\leq6$.

\paragraph{Finite rational certificates.}
All terminating decimals below denote exact rational numbers. Set
\[
 (\ell_1,u_1)=(2.4048,2.4049),\qquad
 (\ell_2,u_2)=(5.5200,5.5201),\qquad c=47/50.
\]
The following polynomial bounds hold throughout the indicated intervals:
\begin{center}
\begin{tabular}{c|c|c}
Polynomial & Interval for $q$ & Bound\\ \hline
$h(q)$ & $[0,9/4]$ & $h(q)\geq1/5$\\
$p(q)$ & $[9/4,\ell_2^2/4]$ & $p(q)\leq-1/50000$\\
$h(q)$ & $[\ell_2^2/4,u_2^2/4]$ & $h(q)\leq-3/25$\\
$p(q)$ & $[(c\ell_1)^2/4,(cu_1)^2/4]$ & $p(q)\leq3081/40000$\\
$p(q)$ & $[(c\ell_2)^2/4,(cu_2)^2/4]$ & $p(q)\leq-11381/100000$
\end{tabular}
\end{center}
For completeness, these are exact coefficient checks in the Bernstein basis.
If $P(q)=\sum_{r=0}^{19}a_rq^r$ on $[u,v]$, define
\begin{equation}\label{bernstein}
 \beta_i=\sum_{k=0}^{i}\frac{\binom{i}{k}}{\binom{19}{k}}
 (v-u)^k\sum_{r=k}^{19}a_r\binom{r}{k}u^{r-k}.
\end{equation}
Then
\[
 P(u+(v-u)t)=\sum_{i=0}^{19}\beta_i\binom{19}{i}t^i(1-t)^{19-i}.
\]
For $0\leq t\leq1$ the basis functions are nonnegative and sum to one.
Substitution of the displayed coefficients of $p$ and $h$ into
\eqref{bernstein} gives the five bounds in the table by comparing each
$\beta_i$ to its stated bound. All coefficient comparisons and polynomial
identities are checked in the Lean certificate; no sampling is used.

\paragraph{Identification of the first two zeros.}
Direct rational evaluation of $p$, together with \eqref{tails}, gives
\begin{equation}\label{signs}
 \J(\ell_1)>0,\quad \J(u_1)<0,\quad
 \J(\ell_2)<0,\quad \J(u_2)>0.
\end{equation}
The first table row and \eqref{tails} imply $H(q)>0$ on $[0,9/4]$.
Consequently $\J$ is strictly decreasing on $[0,3]$, and \eqref{signs}
identifies its unique zero there in $[\ell_1,u_1]$.
The second row implies $\J(x)<0$ for $3\leq x\leq\ell_2$, so there are
no further zeros in that interval. The third row implies $H(x^2/4)<0$ on
$[\ell_2,u_2]$, where $\J$ is strictly increasing. Its endpoint signs give
exactly one zero there. Thus
\begin{equation}\label{zero-bounds}
 j_1\in[\ell_1,u_1],\qquad j_2\in[\ell_2,u_2].
\end{equation}
This proves the zero indices as well as their enclosures.

\begin{proof}[Proof of Theorem~\ref{main}]
Positive fifth powers are strictly increasing. The exact rational inequalities
\[
 1.4204^5<\ell_1^2,\quad u_1^2<1.4205^5,\qquad
 1.9805^5<\ell_2^2,\quad u_2^2<1.9806^5
\]
and $(j_i^{2/5})^5=j_i^2$ give
\[
 1.4204\leq j_1^{2/5}\leq1.4205,\qquad
 1.9805\leq j_2^{2/5}\leq1.9806.
\]
The last two table rows, \eqref{zero-bounds}, and \eqref{tails} yield
\[
 \J(cj_1)\leq\frac{7703}{100000},\qquad
 \J(cj_2)\leq-\frac{569}{5000}.
\]
Using positivity of the power factors and the negative sign of the second
upper bound, we obtain
\begin{align*}
 K_2^{2/5}(c)
 &\leq\frac23\frac{14205}{10000}\frac{7703}{100000}
       -\frac13\frac{19805}{10000}\frac{569}{5000}\\
 &=-\frac{653867}{300000000}<-\frac1{500}.
\end{align*}
\end{proof}

\paragraph{Formal scope and status.}
The certificate contains 87 declarations in eight modules. It proves actual
infinite-series convergence, the derivative identity, uniform tails,
whole-interval polynomial bounds, the first-two-positive-zero identification,
and the literal real-power kernel inequality. The closed endpoint is
\texttt{BesselFejer.source\_counterexample\_365454}.
A fresh isolated parent replay passes with Lean~4.34.1 and mathlib revision
\[\texttt{d13f23b723b8a846827a245b89c10fc7d3f11612},\]
using only \texttt{propext}, \texttt{Classical.choice}, and \texttt{Quot.sound}.
Independent final source, mathematical, formal, manuscript, bounded novelty
and substantive-eligibility review passes. Acceptance additionally requires a
fresh replay of the extracted proof archive. No worldwide priority is claimed.
This is one proposed conjecture refutation; the separate slow-decay existence
conjectures and conditional consequences in the source are not decided here.

\begin{thebibliography}{2}
\bibitem{Babb26} R.~L.~Acosta Babb,
\emph{On the rate of growth of Bessel--Fourier coefficients for integrable functions},
J. Math. Anal. Appl. \textbf{559} (2026), 130440.
\href{https://doi.org/10.1016/j.jmaa.2026.130440}{doi:10.1016/j.jmaa.2026.130440}.
\bibitem{Babb24} R.~L.~Acosta Babb,
\emph{A partial converse to the Riemann--Lebesgue lemma for Bessel--Fourier
series of order zero}, arXiv:2410.17681v1 (2024).
\end{thebibliography}
\end{document}
